\section{区域秩序不是统一的空白}
\label{sec:five-dynasties-regional-order}

907年哀帝禅位时，开封的军队、汴河和宣武旧部让朱温能够把后梁的国号安在一座可供给的中枢上；同一年，成都、凤翔、淮南、两浙和岭南面对的却不是同一份选择题。王建称帝、李茂贞沿用唐年号、杨氏拒梁命、钱镠接受王号，分别把军府、地方官僚、港口或腹地接到不同的名义语言上。称帝、奉唐或受封都不能直接换算为行政密度；它们首先是让地方能够任命、交涉、征发和安排继承的制度接口。%
\footnote{唐亡与各区域政权的重新定位见 LT-907-080 至 LT-907-085，依据《旧唐书·哀帝纪》《旧五代史》梁太祖纪及王建、李茂贞、杨行密、钱镠诸传，并可对读《新五代史》诸世家。仪式文本、王号和年号不能替代州县税籍或基层服从的记录。}

区域之间也没有一条共同的兴亡时间。成都凭盆地道路、粮区和既有军府承受中原战争的余波；杭州与两浙依靠水网、港口和地方官僚维系吴越；淮南、湖湘、荆江、闽岭与岭南则各有河道、山地、盐粮和城镇关系。923年后唐灭梁、925年灭前蜀，旋即在926年遭兵变，934年后蜀又在成都建立，最能说明战场接收并未消去地方的兵粮与人事网络。北方王朝的胜利可以改变谁占有一座都城，却难以一次改写每个区域的供给、记忆和协商方式。%
\footnote{后唐灭梁、灭前蜀、庄宗遇弑和后蜀建立见 LT-923-092、LT-925-094、LT-926-095、LT-934-098。《旧五代史》《新五代史》与《资治通鉴》提供中原年序；蜀地接收和地方损失仍须结合墓志、城址和区域材料，不能由军功叙事补足。}

北方边界使开封的继承从来不是仅限宫门的事务。936年石敬瑭需要契丹援军以挑战后唐，契丹则要求燕云州郡和政治承认；这笔交易让城、牧地、关隘与互市成为后来军费的一部分。后晋灭亡、辽军北撤、后汉和后周相继进入开封，不能被写成单纯的“中原复失”。每一方都要面对同一个困难：远离自身供给基地的军队怎样获得粮食和地方合作，新的中枢怎样接住旧朝的禁军、官僚和税源。%
\footnote{后晋建立、后晋覆亡、辽军北撤与后汉建立见 LT-936-099、LT-946-105、LT-947-106、LT-947-107，依据《旧五代史》晋、汉纪、《辽史·太宗纪》与《资治通鉴》。燕云割让文书原件不存，范围、礼仪和撤军原因均有材料缺环。}

后周对南唐的战争显示“区域秩序”也会被改造，而非只是被保存。956年沿淮南下，958年接收江北州县后，后周仍须安置降军、修复城池、恢复仓储和水路，才可能把淮河以北从战区变为可供北防的财政空间。955年的寺院整顿同样说明军费、铜料、户籍和地方机构彼此相连。960年赵匡胤受禅继承的是开封中枢和禁军，不是一张已经统一的地图；辽、北汉和南方诸国仍在，宋必须通过后续战争、谈判和接收才逐段处理这些关系。%
\footnote{后周整顿寺院、南征、江北接收与陈桥受禅见 LT-955-112、LT-956-113、LT-958-115、LT-960-118。可据《旧五代史·周世宗纪》《五代会要》《宋史·太祖本纪》及《资治通鉴》；寺院执行、割地范围和陈桥预谋均不能由单一正统叙事裁定。}

\paragraph{多中心世界的制度得失。}
五代十国的区域政权降低了所有资源都必须经开封才能组织的风险：不同城市和军府能够维持贸易、任官、宗教与地方秩序。代价是中原朝廷难以把一次诏令或一次胜利变成稳定的税兵控制，北方边界又不断提高继承与防务的价格。把这段历史写成“统一前的混乱”，会抹掉这些区域秩序的实际工作；把它们写成彼此隔绝的国家，也会看不见水路、册命、市场和战争持续把它们连在一起。
